\subsection{Chern and Pontryagin Classes}

    Let $E \to M$ be a complex vector bundle of rank $d$. For any matrix $A \in \C^{d \times d}$ we can consider the characteristic polynomial in $t$:
    \[
    \mathrm{det}\left( \mathrm{id} - \frac{t}{2\pi i} A\right) = \sum_{k=0}^{d} f_k(A)t^k,
    \]
    where $f_k \colon \C^{d \times d} \to \C$ are the coefficients of this polynomial and depend on $A$. Using the Leibniz formula for the determinant one can see that $f_k$ is a $k$-homogeneous polynomial and it moreover is invariant since the characteristic polynomial is. Furthermore it follows from the Leibniz formula that 
    \begin{align*}
        f_0(A) & = 1, \\
        f_1(A) &= \mathrm{tr}\left(-\frac{1}{2\pi i} A\right), \\
        f_d(A) &= \mathrm{det}\left(-\frac{1}{2\pi i} A\right).
    \end{align*}

\begin{definition}
    The \textbf{$k$-th Chern class} of $E$ is defined as 
    \[
    c_k(E) := \chi_E(f_k) \in H_\mathrm{dR}^{2k}(M;\C).
    \]
\end{definition}

If we instead consider a real vector bundle $E \to M$ we can similiarly define $k$-homogeneous polynomials $g_k \colon \R^{d \times d} \to \R$ by 
\[
    \mathrm{det}\left( \mathrm{id} - \frac{t}{2\pi} A\right) = \sum_{k=0}^{d} g_k(A)t^k,
\]
and define:

\begin{definition}
    The \textbf{$k$-th Pontryagin class} of $E$ is defined as 
    \[
    p_k(E) := \chi_E(g_{2k}) \in H_\mathrm{dR}^{4k}(M;\R).
    \]
\end{definition}

Of course there is a connection between these two definitions:

\begin{lemma}
    Let $E \to M$ be a real vector bundle. Then 
    \[
    p_{k}(E) = (-1)^k c_{2k}(E \otimes \C),
    \]
    where $E \otimes \C$ is the complexified bundle and the equality is understood by considering a real differential as the real part of a complex differential form, that is $H_\mathrm{dR}^{4k}(M;\R) \subseteq H_\mathrm{dR}^{4k}(M;\C)$.
\end{lemma}

\begin{proof}
    Choose any connection $\nabla$ on $E$ and let $\nabla_\C$ be the induced connection on the complexified bundle $E \otimes \C$. Then by definition any frame of $E$ also defines a frame of $E \otimes \C$ and the corresponding connection $1$-forms are equal. Thus the curvature form $\Omega$ of $E$ agrees with the curvature form $\Omega_\C$ of $E \otimes \C$. From the definition of the polynomials $f_k$ and $g_k$ it is moreover clear that for any real matrix $A \in \R^{d \times d}$ we have 
    \[
    f_k(A) = (-i)^k g_k(A).
    \]
    Thus 
    \[
    c_{2k}(E) = [f_{2k}(\Omega_\C)] = [f_{2k}(\Omega)] = (-i)^{2k} [g_{2k}(\Omega)] = (-1)^k p_k(E).
    \]
\end{proof}

Now assume we are given a real oriented vector bundle $E \to M$ of rank $2n$. Choose a metric on $E$ together with a compatible connection. A simple computation shows that the connection $1$-form $\Theta$ with respect to any oriented orthonormal frame is skew-symmetric. Thus 
\begin{align*}
    (\Theta \wedge \Theta)^T(X,Y) &= (\Theta(X)\Theta(Y) - \Theta(Y)\Theta(X))^T \\ 
    &= (-\Theta(Y))(-\Theta(X)) - (-\Theta(X))(-\Theta(Y)) \\
    &= - (\Theta \wedge \Theta) (X,Y).
\end{align*}
So the curvature form $\Omega$ is also skew symmetric:
\[
\Omega^T = (d\Theta)^T + (\Theta \wedge \Theta)^T = -d\Theta - \Theta \wedge \Theta = -\Omega.
\]
Recall that the Pfaffian is a $n$-homogeneous polynomial 
\[
\mathrm{Pf} \colon \mathrm{Skew}_{2n} \to \R
\]
on the space of skew symmetric $2n \times 2n$ matrices and satisifies 
\[
\mathrm{Pf}(B^T A B) = \mathrm{det}(B)\mathrm{Pf}(A)
\]
for $A \in \mathrm{Skew}_{2n}$ and $B \in \R^{n \times n}$ an arbitrary matrix. Thus we can define the $2n$-form $\mathrm{Pf}(\Omega)$ on the neighbourhood where the oriented orthonormal frame was defined. Under a change of oriented orthonormal frame the curvature form $\Omega$ transforms as $S \Omega S^{-1}$ for a special orthogonal matrix $S \in \mathrm{SO}(2n)$. Since $S^{-1} = S^T$ and $\mathrm{det}(S) = 1$ we conclude that $\mathrm{Pf}(\Omega)$ does not depend on the chosen frame and thus defines a global object. Just as in the proof of Theorem \ref{chern_weil:thm:chern_weil_thm} we see that $\mathrm{Pf}(\Omega)$ is closed (only the statement of Lemma \ref{chern_weil:lem:technical_lemma1} changes, where instead of arbitrary $X_i$ and $Y$ we only allow for skew symmetric matrices. Note that the tangent space of $\mathrm{SO}(2n)$ is precisely the space of skew symmetric matrices, such that we can realize any such $Y$ as the velocity vector of a curve through $\mathrm{id}$). Moreover the cohomology class does not depend on the chosen connection and metric. Again the proof is similiar to Theorem \ref{chern_weil:thm:chern_weil_thm}. Thus we can define 

\begin{definition}
    The \textbf{Euler class} of $E$ is defined as 
    \[
    e(E) := \left[\mathrm{Pf}\left(\frac{1}{2\pi}\Omega\right)\right] \in H_\mathrm{dR}^{2n}(M;\R).
    \]
\end{definition}

It is no coincidence that the above construction of the Euler class works and is just an example of the following more general construction. Instead of considering homogeneous polynomials on $\K^{k \times k}$ which are invariant under conjugation with elements in $\mathrm{GL}(k,\K)$ we can more generally consider polynomials $\eta \colon\mathfrak{g} \to \K$ defined on the Lie-Algebra $\mathfrak{g}$ of some matrix Lie-group $G \subseteq \mathrm{GL}(k,\K)$, which are invariant under conjugation with elements in $G$. Suppose we are then given a vector bundle $E \to M$ with structure group $G$ and some connection, such that the corresponding (locally defined) curvature form $\Omega$ takes values in $\mathfrak{g}$. Under a change of local frame this curvature form transforms as $S \Omega S^{-1}$ for a matrix $S \in G$, thus $\eta(\Omega)$ is a well defined global object and the proofs above are easily generalized to show that $\eta(\Omega)$ is closed and does depend on the chosen connection (as long as the connection yields local curvature forms with values in $\mathfrak{g}$). For this more general approach to be useful we need to guarantee the existence of such connections, which is the statement of the following proposition.

\begin{proposition}
    Let $E \to M$ be a vector bundle with structure group $G \subseteq \mathrm{GL}(k,\K)$. Then there exists a connection on $E$, such that the curvature form with respect to any local frame \footnotemark takes values in $\mathfrak{g}$. 
\end{proposition}

\begin{proof}
    Let $(e_i^\alpha)$ be local frames defined on some neighbourhoods $U_\alpha \subseteq M$ which cover $M$. On $E|_{U_\alpha}$ define a connection by $\nabla^\alpha e_i = 0$. In other words for a local section $s = s^i e_i^\alpha$ we set $\nabla^\alpha s := d s^i \otimes e_i^\alpha$. Now let $(\psi_\alpha)$ be a partition of unity subordinate to the cover $(U_\alpha)$ and define a connection $\nabla$ on $E$ by 
    \[
    \nabla s := \sum_{\alpha} \psi_\alpha \nabla^\alpha s
    \]
    for any section $s \in \Gamma(E)$. We claim that the connection one forms of $\nabla$ already take values in $\mathfrak{g}$. To this end let $(e_i)$ be any other frame defined on $U \subseteq M$ and let $A_\alpha = (a_{ij}^\alpha) \in C^\infty(U \cap U_\alpha,G)$ be the change of basis matrix from $(e_i)$ to $(e_i^\alpha)$, that is 
    \[
    e_j = \sum_{i} a_{ij}^\alpha e_i^\alpha.
    \]
    Moreover we write $A_\alpha^{-1} = (a^{ij}_\alpha)$ for the inverse, that is 
    \[
    e_j^\alpha = \sum_{i} a^{ij}_\alpha e_i.
    \]
    We now compute 
    \begin{align*}
        \nabla^\alpha e_j = \sum_i d a_{ij}^\alpha \otimes e_i^\alpha = \sum_{i,k} a_\alpha^{ki} da_{ij}^\alpha \otimes e_k = \sum_{k,i} a_\alpha^{ik} da_{kj}^\alpha \otimes e_i 
    \end{align*}
    In the last equality we only exchanged the roles of $k$ and $i$. So 
    \[
    \nabla e_j = \sum_{\alpha,k,i} \psi_\alpha a_\alpha^{ik} da_{kj}^\alpha \otimes e_i,
    \]
    such that for some vector field $X$ the connection one form is given by 
    \[
    \Theta(X) = \sum_\alpha \psi_\alpha \left(\sum_k a_\alpha^{ik} da_{ki}^\alpha(X) \right)_{ij} = \sum_\alpha \psi_\alpha A_\alpha^{-1} dA_\alpha(X)
    \]
    Note that for some point $p \in M$, where both frames $(e_i)$ and $(e_i^\alpha)$ are defined we have $$d_pA_\alpha(X_p) \in T_{A_\alpha(p)} G,$$ such that $A_\alpha^{-1}(p) d_pA_\alpha(X_p) = d_p(A_\alpha^{-1}(p)A_\alpha)(X_p) \in T_e G = \mathfrak{g}$. Thus $\Theta(X) \in \mathfrak{g}$. Next we show that the corresponding curvature form 
    \[
    \Omega = \Theta \wedge \Theta + d\Theta
    \]
    also takes values in $\mathfrak{g}$. Firstly note that for vector fields $X,Y$ we have 
    \[  
    (\Theta \wedge \Theta)(X,Y) = \Theta(X)\Theta(Y) - \Theta(Y)\Theta(X) = [\Theta(X),\Theta(Y)]
    \]
    and it follows from general Lie theory of matrix groups, that the corresponding Lie algebra is closed under taking the commutator. For the second term we apply Corollary \ref{cartan_cal:cor:formula_d_alpha} component wise to get 
    \[
    d\Theta(X,Y) = X \cdot \Theta(Y) - Y \cdot \Theta(X) - \Theta([X,Y]).
    \]
    The first two terms are just directional derivatives of smooth functions with values in the vector space $\mathfrak{g}$ and thus take again values in $\mathfrak{g}$. The statement follows.
\end{proof}

\footnotetext{That is any local frame determined by a local trivialisation from the maximal trivialising atlas making $E \to M$ into a vector bundle with structure group $G$. Compare Definition \ref{fiber_bundles:def:fiber_bundle}.}